\documentclass[11pt,a4paper]{article}
\usepackage{fontspec}
\usepackage[finnish]{babel}
\usepackage[a4paper,margin=23mm,headheight=15pt,footskip=18mm]{geometry}
\usepackage{amsmath,amssymb,mathtools}
\usepackage{enumitem}
\usepackage{xcolor}
\usepackage{xurl}
\usepackage{hyperref}
\usepackage{fancyhdr}
\usepackage{titlesec}
\usepackage{microtype}
\usepackage{array}
\usepackage{bookmark}

\setmainfont{DejaVu Serif}
\setsansfont{DejaVu Sans}
\setmonofont{DejaVu Sans Mono}

\definecolor{TitleBlue}{HTML}{17233A}
\definecolor{AccentBlue}{HTML}{355F96}
\definecolor{PaleBlue}{HTML}{EEF4FA}
\hypersetup{
  unicode=true,
  colorlinks=true,
  linkcolor=AccentBlue,
  urlcolor=AccentBlue,
  pdftitle={Puolijakson symmetriat modulo 3},
  pdfsubject={Matemaattinen formalisointi ja algoritminen auditointi},
  pdfkeywords={recurrence, modulo 3, antiperiodicity, Fourier, companion matrix},
  pdfcreator={XeLaTeX}
}

\titleformat{\section}{\large\bfseries\color{TitleBlue}}{\thesection.}{0.55em}{}
\titleformat{\subsection}{\normalsize\bfseries\color{TitleBlue}}{\thesubsection}{0.55em}{}
\titlespacing*{\section}{0pt}{1.2ex plus .2ex}{.55ex}
\titlespacing*{\subsection}{0pt}{.9ex plus .2ex}{.35ex}
\setlist[itemize]{leftmargin=1.7em,itemsep=.20em,topsep=.35em}
\setlist[enumerate]{leftmargin=1.9em,itemsep=.20em,topsep=.35em}
\setlength{\parindent}{0pt}
\setlength{\parskip}{.48em}
\emergencystretch=2.5em
\setcounter{tocdepth}{2}

\pagestyle{fancy}
\fancyhf{}
\fancyhead[L]{\small\color{AccentBlue}Puolijakson symmetriat modulo 3}
\fancyhead[R]{\small\color{AccentBlue}Formalisointi ja auditointi}
\fancyfoot[C]{\small\color{TitleBlue}\thepage/6}
\renewcommand{\headrulewidth}{0.35pt}
\renewcommand{\footrulewidth}{0pt}

\newcommand{\F}{\mathbb{F}}
\newcommand{\Z}{\mathbb{Z}}
\newcommand{\oplusword}{\mathbin{\Vert}}
\newcommand{\centerlift}[1]{\widetilde{#1}}

\begin{document}
\sloppy
\thispagestyle{empty}
\vspace*{20mm}
\noindent\noindent{\color{AccentBlue}\rule{\textwidth}{0.8pt}}\\[22mm]
{\raggedright\fontsize{24}{29}\selectfont\bfseries\color{TitleBlue}Puolijakson symmetriat\\ modulo 3\par}
\vspace{9mm}
{\fontsize{15}{20}\selectfont\color{AccentBlue}Matemaattinen formalisointi ja algoritminen auditointi\par}
\vspace{17mm}
\noindent\fcolorbox{AccentBlue}{PaleBlue}{%
  \begin{minipage}{0.92\textwidth}
  \vspace{2mm}\small Tässä muistiossa erotetaan täsmällisesti kolme usein toisiinsa sekoitettua ilmiötä: kiertojen negaatio, puolijakson antiperiodisuus ja kääntö. Siinä todistetaan nollasta poikkeavan primitiivisen sanan yksikäsitteinen vastakkaissiirto, annetaan kombinatoriset, spektraaliset ja matriisimuotoiset karakterisoinnit sekä liitetään tulokset modulaaristen rekursioiden kattavaan analysointiin. \vspace{2mm}
  \end{minipage}}
\vfill
{\color{TitleBlue}\small
Tietovarasto: \textsf{Modular-recurrence-symetry-analyzer-}\\
Yhdistetty PDF-versio - 5. elokuuta 2026
}\par
\vspace{18mm}
\noindent{\color{AccentBlue}\rule{\textwidth}{0.8pt}}
\clearpage
\setcounter{page}{1}

\tableofcontents
\clearpage


\section{Tarkoitus}
Tässä muistiossa formalisoidaan ilmiöt, joita tietovarastossa \nolinkurl{59200/k-bonacci-pisano-finder} kuvataan epämuodollisilla ilmauksilla ``kolmea täydentävät jaksot'' ja ``itseään täydentävä kahtia leikattaessa''.

Invariantti ympäristö on jaksollisten sanojen joukko yli kunnan
\[
\F_3=\{0,1,2\},
\]
varustettuna additiivisella involuutiolla
\[
\nu(x)=-x\pmod3,\qquad \nu(0)=0,\quad \nu(1)=2,\quad \nu(2)=1.
\]
Ilmaus ``3:n komplementti'' tulee siksi korvata käsitteellä \emph{additiivinen käänteisluku modulo 3} tai \emph{negaatiokomplementti modulo 3}. Valittujen kokonaislukuedustajien desimaalinen summa ei ole aina 3, koska 0 kuvautuu luvuksi 0.

\section{Kolme erillistä ominaisuutta}
Olkoon $w=(w_0,\ldots,w_{T-1})\in\F_3^T$ sana, jonka primitiivinen jakso on $T$.

\subsection{Vastakkaisten jaksojen pari}
Kaksi jaksoa $w$ ja $v$ muodostavat vastakkaisen parin, jos
\[
v_n=-w_n\pmod3
\]
kaikilla $n$. Ne voivat kuulua kahteen eri kiertoon.

\subsection{Puolijakson antiperiodisuus}
Sana $w$ on puolijakson suhteen antiperiodinen, jos $T=2h$ ja
\[
w_{n+h}=-w_n\pmod3
\]
kaikilla $n$. Tällöin se voidaan kirjoittaa muodossa
\[
w=H\oplusword(-H),\qquad H=(w_0,\ldots,w_{h-1}).
\]
Esimerkkejä:
\[
01120221=0112\oplusword0221,\quad 011022=011\oplusword022,\quad 01220211=0122\oplusword0211.
\]

\subsection{Jälkimmäisen puoliskon kääntö}
Operaatio
\[
R(a_0,\ldots,a_{h-1})=(a_{h-1},\ldots,a_0)
\]
on riippumaton negaatiosta. Fibonacci-esimerkissä
\[
R(0221)=1220.
\]
Niinpä 1220 on sanan 0112 käännetty antipodaalinen sana, ei jälkimmäinen puolisko sellaisenaan.

\clearpage
\section{Yksikäsitteisen vastakkaissiirron lause}
\textbf{Lause.} Olkoon $w$ nollasta poikkeava primitiivinen sana, jonka jakso on $T$. Jos on olemassa siirto $r$ siten, että
\[
w_{n+r}=-w_n
\]
kaikilla $n$, niin $T$ on parillinen ja
\[
r\equiv\frac{T}{2}\pmod T.
\]

\textbf{Todistus.} Kun siirtoa sovelletaan kahdesti, saadaan $w_{n+2r}=w_n$. Primitiivisyys antaa ehdon $T\mid2r$. Siirto $r$ ei ole nolla modulo $T$, koska nollasta poikkeava sana yli $\F_3$:n ei voi toteuttaa yhtälöä $w=-w$. Syklisen ryhmän $\Z/T\Z$ ainoa epätriviaali kertalukua 2 oleva alkio on olemassa vain, kun $T$ on parillinen, ja tällöin se on $T/2$. \hfill$\square$

\textbf{Seuraus.} Nollasta poikkeava primitiivinen jakso on siis joko
\begin{enumerate}
\item negaation eri kiertoon kuvaama; tai
\item negaation suhteen invariantti, jolloin se on välttämättä puolijakson antiperiodinen.
\end{enumerate}
Tämä kahtiajako poistaa sekaannuksen parin
\[
00111201,\qquad 00222102
\]
ja sisäisesti antiperiodisten jaksojen, kuten 011022, välillä.

\section{Kombinatoriset seuraukset}
Jos $w=H\oplusword(-H)$, niin
\[
\#\{n:w_n=1\}=\#\{n:w_n=2\},\qquad \#\{n:w_n=0\}\equiv0\pmod2,
\]
ja
\[
\sum_{n=0}^{T-1}w_n=0\quad\text{kunnassa }\F_3.
\]
Ehdot ovat välttämättömiä mutta eivät riittäviä. Keskitetyn noston
\[
\centerlift0=0,\qquad \centerlift1=1,\qquad \centerlift2=-1
\]
avulla generoiva polynomi tekijöityy muodossa
\[
W(X)=\sum_{n=0}^{2h-1}\centerlift{w_n}X^n=(1-X^h)\sum_{n=0}^{h-1}\centerlift{w_n}X^n.
\]

\clearpage
\section{Fourier-karakterisointi}
Olkoon
\[
\widehat w(k)=\sum_{n=0}^{2h-1}\centerlift{w_n}e^{-2\pi i kn/(2h)}.
\]
Ehto $w_{n+h}=-w_n$ pätee täsmälleen silloin, kun $\widehat w(k)=0$ jokaisella parillisella indeksillä $k$. Kun $w_{n+h}=-w_n$,
\[
\widehat w(k)=\bigl(1-(-1)^k\bigr)\sum_{n=0}^{h-1}\centerlift{w_n}e^{-2\pi i kn/(2h)},
\]
joten kerroin häviää parillisilla $k$:n arvoilla. Kääntäen, jos kaikki parilliset moodit häviävät, sana kuuluu parittomien moodien virittämään aliavaruuteen. Siirto $h$:lla vaikuttaa moodiin $k$ kertoimella $(-1)^k=-1$ ja siis sanaan operaattorina $-I$. Kyseessä on täsmällinen spektraalinen testi, ei pelkkä kahden puoliskon visuaalinen vertailu.

\section{Lineaariset jonot modulo 3}
Tarkastellaan kertaluvun $k$ rekursiota
\[
u_{n+k}=c_0u_n+c_1u_{n+1}+\cdots+c_{k-1}u_{n+k-1}\pmod3.
\]
Tila on $s_n=(u_n,\ldots,u_{n+k-1})^T$, ja kehitys on $s_{n+1}=Ms_n$, missä $M$ on kumppanimatriisi.

\subsection{Puhdas jaksollisuus}
Matriisin $M$ determinantti on etumerkkiä lukuun ottamatta $c_0$. Kunnassa $\F_3$
\[
M\text{ on kääntyvä}\iff c_0\neq0.
\]
Tällöin jokainen tilarata on puhtaasti jaksollinen. Jos $c_0=0$, voi esiintyä esijakso, joka on erotettava varsinaisesta kierrosta.

\subsection{Kiertojen negaatio}
Lineaarisuus antaa $M(-s)=-M(s)$. Negaatio kuvaa siis jokaisen kierron yhtä pitkään kiertoon. Nollasta poikkeavalle kierrolle on vain kaksi mahdollisuutta:
\begin{itemize}
\item kaksi eri kiertoa vaihtuvat negaation vaikutuksesta;
\item yksi invariantti kierto, joka on välttämättä puolijakson antiperiodinen.
\end{itemize}
Nollakierto on ainoa degeneroitunut poikkeus: negaatio kiinnittää sen pisteittäin, eikä se määritä epätriviaalia primitiivistä antijaksoa.

\subsection{Yhden alkutilan ehto}
Alkutilalle $s_0$ relaatio
\[
M^hs_0=-s_0
\]
implikoi jokaiselle tilan lineaariselle ulostulolle
\[
u_{n+h}=-u_n.
\]
Kun $s_0\neq0$ ja kierto on primitiivinen, relaatio tuottaa edellä kuvatun puolijakson antijakson.

\clearpage
\subsection{Yhtenäinen ehto}
Kaikki alkutilat toteuttavat saman antijaksorelaation $h$ täsmälleen silloin, kun
\[
M^h=-I.
\]
Yhtäpitävästi minimipolynomi $\mu_M$ toteuttaa
\[
\mu_M(X)\mid X^h+1\qquad\text{kunnassa }\F_3[X],
\]
ja siten $M^{2h}=I$.

\section{Fibonaccin esimerkki}
Kun
\[
M=\begin{pmatrix}0&1\\1&1\end{pmatrix}\quad\text{kunnassa }\F_3,
\]
saadaan $M^4=-I$ ja $M^8=I$. Alkutilan $(0,1)$ generoima jakso on
\[
01120221=0112\oplusword(-0112).
\]
Käännetty jälkimmäinen puolisko on $R(0221)=1220$. Desimaalilukuna luettaessa havaittu numeerinen yhtälö
\[
1220=L_{14}+F_{14}=2F_{15}
\]
on ylimääräinen koodausidentiteetti eikä antiperiodisuuden yleinen seuraus.

\section{Yleistys mielivaltaisiin jonoihin modulo 3}
Annetun jaksollisen sanan analysointi ei edellytä rekursio-oletusta. Jokaiselle primitiiviselle jaksolle voidaan laskea
\begin{itemize}
\item sen rotaatiokierto;
\item sen additiivinen käänteisluku;
\item sen affiiniset symmetriat $w_{n+r}=aw_n+b$;
\item sen kääntösymmetriat $w_{r-n}=aw_n+b$;
\item mahdollinen antiperiodisuus;
\item sen parilliset ja parittomat Fourier-moodit.
\end{itemize}
Matriisiteoriaa tarvitaan vain, kun annettu lineaarinen rekursio on käytettävissä.

\section{Laajennus yleiseen modulukseen}
Renkaassa $\Z/m\Z$ luonnollinen involuutio on edelleen $x\mapsto-x$. Määritelmä
\[
w_{n+h}=-w_n\pmod m
\]
säilyy voimassa. Kun $m=2$, negaatio on identiteetti ja antiperiodisuus supistuu tavalliseksi jaksollisuudeksi. Kun $m>2$, ero on epätriviaali. Ilmausta ``$m$:n komplementti'' ei pidä käyttää algebrallisena määritelmänä, koska se riippuu valituista kokonaislukuedustajista; invariantti termi on ``additiivinen käänteisluku modulo $m$''.

\clearpage
\section{Algoritminen auditointi}
Kertaluvun $k$ rekursion jakso modulo $m$ on havaittava äärellisessä tila-avaruudessa
\[
(\Z/m\Z)^k,
\]
ei etsimällä toistuvaa lohkoa mielivaltaisen pitkästä prefiksistä. Toimitettu ohjelma käyttää
\begin{itemize}
\item kompaktia kanta-$m$-tilakoodausta;
\item Brentin algoritmia yhdelle alkutilalle, ajassa $O(\mu+\lambda)$ ja muistissa $O(1)$;
\item kaikkien tilojen funktionaalisen graafin täsmällistä hajotelmaa ajassa $O(m^k)$;
\item rinnakkaistamista kerroinperheiden yli;
\item vastakkaisten kiertojen ja antiperiodisten kiertojen erillistä luokittelua;
\item matriisitarkistusta $M^h=-I$.
\end{itemize}
Mathematica-ydintä tai ulkoista käärettä ei tarvita.

\section{Toistettavat skannaustulokset}
Tietovaraston viidentoista nimetyn perheen skannaus löytää tritetranacci-perheestä erilliset vastakkaiset parit, joiden pituudet ovat 24 ja 8, sekä antiperiodiset jaksot 011022, 01220211 ja 12.

Toinen skannaus kattaa 119 nollasta poikkeavaa binääristä kerroinvektoria kertaluvuilla 2--6:
\[
\sum_{k=2}^{6}(2^k-1)=119.
\]
Ohjelma löytää
\begin{itemize}
\item 13 perhettä, jotka täyttävät globaalin ehdon $M^h=-I$;
\item 88 perhettä, joissa on vähintään yksi epätriviaali puolijakson antiperiodinen kierto;
\item 96 perhettä, joissa on vähintään yksi erillinen vastakkaisten kiertojen pari.
\end{itemize}
Tulokset ovat tyhjentäviä tässä äärellisessä tarkasteluikkunassa. Ne voidaan toistaa tiedostoista \nolinkurl{mod3_symmetry_scan_summary.json}, \nolinkurl{binary_recurrence_families_mod3.json} ja \nolinkurl{binary_recurrence_families_mod3_summary.csv}.

\section*{Johtopäätös}
Negaatio modulo 3 toimii rakenteellisena involuutiona jaksollisissa sanoissa ja lineaaristen rekursioiden kierroissa. Nollasta poikkeavan primitiivisen sanan invarianssi tämän involuution suhteen voi toteutua vain yksikäsitteisen puolijaksosiirron kautta. Sama ominaisuus tunnistetaan yhtäpitävästi generoivan polynomin tekijöitymisestä, parillisten Fourier-moodien häviämisestä tai lineaarisessa asetelmassa matriisirelaatiosta $M^h=-I$. Algoritminen analyysi erottaa siten kiertojen sisäiset symmetriat pelkistä vastakkaisten kiertojen välisistä yhteyksistä.


\end{document}
